\documentclass[10pt,a4paper]{amsart}
\usepackage[margin=19mm]{geometry}
\usepackage{amsmath,amssymb,mathtools}
\usepackage[scr=rsfs,cal=euler]{mathalfa}
\usepackage{xfrac}
\usepackage[unicode,hidelinks,bookmarks=false]{hyperref}
\newcommand{\ie}{i.e.\ }
\newcommand{\cf}{cf.\ }
\newcommand{\resp}{respectively\ }
\newcommand{\Subsection}{\subsection}
\providecommand{\nobreakdash}{\nobreak\hbox{-}}
\setlength{\emergencystretch}{2em}
\multlinegap0pt
\usepackage{enumerate}
\usepackage[shortlabels]{enumitem}
\usepackage{amssymb}
\usepackage{bm}
\usepackage{cancel}
\newtheorem{assumption}{Assumption}
\newtheorem{corollary}{Corollary}
\newtheorem{proposition}{Proposition}
\title{Contraction of hypersurfaces with positive sectional curvature in hyperbolic space}
\author{Tianci Luo, Yong Wei, Rong Zhou}
\date{Source: arXiv:2604.25513v1 (CC BY 4.0)}
\begin{document}
\maketitle

\noindent\emph{Source and licence.} Contraction of hypersurfaces with positive sectional curvature in hyperbolic space. Version arXiv:2604.25513v1 (CC BY 4.0), distributed by arXiv under the Creative Commons Attribution 4.0 International licence, http://creativecommons.org/licenses/by/4.0/, as recorded in the arXiv licence metadata for this exact version and shown on the abstract page https://arxiv.org/abs/2604.25513v1. Full licence notice: \texttt{licenses/arxiv-2604.25513-CC-BY-4.0.txt}.\par\smallskip

\noindent\textit{Editorial scope.} Counted unit: the article's proposition s4.lem-pinc (source0.tex lines 487--494). The article's complete original proof of it is delivered with it (source0.tex lines 495--662, one \texttt{proof} environment of 168 lines). No mathematics is written, repaired or re-derived: the statement and the proof are the article's own lines, in the article's own notation, and no retained line is substituted or re-flowed.  Retained uncounted as the source-local context the proof needs: lines 81--89; lines 114--135; lines 205--208; lines 261--267; lines 380--386; lines 382--385; lines 409--415; lines 411--414.  Nothing was omitted from any retained span, so the package carries no omission record.  No retained line contains a citation, so the wrapper carries no bibliography and no bibliography entry.  No retained line contains a figure environment, an image-inclusion command or drawing code, so no image is delivered and the package declares no asset dependency.  Editorial formatting only, in addition: an AMS \texttt{amsart} preamble that restates the article's own declarations and macros used by the retained lines, namely the environments \texttt{assumption}, \texttt{corollary}, \texttt{proposition} on separate counters, together with the standard packages those lines need.  The retained environments print in this document as Assumption 1, Corollary 1, Proposition 1 in order of appearance, whereas the article prints the same items with the numbering of its own full document.  The complete native source of the article as distributed in the arXiv e-print for this exact version is bundled unchanged as \texttt{sources/199309/source0.tex}.
\begin{equation}
	\label{equ-f}
	\begin{cases}
		\begin{aligned}
			\dfrac{\partial}{\partial t}X(x,t) & = -F(x,t)\,\nu(x,t), \\
			X(\cdot,0) & = X_0,
		\end{aligned}
	\end{cases}
\end{equation}

\begin{assumption}	\label{assumption}
$F:\Gamma_+\to (0,\infty)$ is a smooth symmetric function defined on the positive cone
	\begin{equation*}
		\Gamma_+ := \{ \kappa=(\kappa_1,\kappa_2,\cdots,\kappa_n)\in\mathbb{R}^n:\kappa_i>0,\ \forall\ i=1,2,\cdots,n \},
	\end{equation*}
	and satisfies the following properties:
	\begin{enumerate}
		\item[{\rm (i)}] $F$ is strictly increasing in each variable, homogeneous of degree one, and normalized by $F(1,1,\cdots,1)=1$;
		%\item[{\rm (ii)}] For any $i\neq j$,
		%\begin{equation}
		%	\left( \dfrac{\partial F}{\partial \kappa_i} \kappa_i - \dfrac{\partial F}{\partial \kappa_j} \kappa_j  \right)(\kappa_i - \kappa_j) \geq 0;
		%	\label{equ-pp2}
		%\end{equation}
		\item[{\rm (ii)}] For all $(y_1,y_2,\dots,y_n)\in\mathbb{R}^n$,
		\begin{equation}
			\sum\limits_{k,\ell}\ddot{F}^{k\ell}y_ky_{\ell} \geq F^{-1} \left(\sum\limits_{k=1}^n \dot{F}^k y_k \right)^2 - \sum\limits_{k=1}^n \dfrac{\dot{F}^k}{\kappa_k}y_k^2,
			\label{equ-property3}
		\end{equation}
        where $\dot{F}^k=\frac{\partial F}{\partial \kappa_k}$ and $\ddot{F}^{k\ell}=\frac{\partial^2F}{\partial\kappa_k\partial\kappa_\ell}$. 
		%\item[{\rm (iv)}] $f$ is concave.
	\end{enumerate}
\end{assumption}

\begin{equation}
	\ddot{F}^{ij,k\ell}(A)B_{ij}B_{k\ell} = \sum\limits_{i,j}\ddot{F}^{ij}(\kappa)B_{ii}B_{jj} + 2 \sum\limits_{i>j} \dfrac{\dot{F}^i(\kappa)-\dot{F}^j(\kappa)}{\kappa_i-\kappa_j}B_{ij}^2.
	\label{equ-secf}
\end{equation}

\begin{align}
    \frac{\partial}{\partial t}g_{ij}=&-2Fh_{ij},\label{s2.evl-g}\\
		\frac{\partial}{\partial t} F = & \dot{F}^{k\ell}\nabla_k\nabla_{\ell}F + \left( \dot{F}^{ij}(h^2)_{ij}- \dot{F}^{ij}g_{ij} \right)F,  \label{equ-psi}\\
        \frac{\partial}{\partial t}h_{ij}=&\nabla_i\nabla_jF-F\left((h^2)_{ij}+g_{ij}\right),\label{s2.evl-h1}\\
        \frac{\partial}{\partial t}h_{ij}=&\dot{F}^{k\ell}\nabla_k\nabla_{\ell}h_{ij} + \ddot{F}^{k\ell,pq}\nabla_i h_{k\ell} \nabla_j h_{pq} \nonumber\\
        &\quad + \left( \dot{F}^{k\ell}(h^2)_{k\ell} + \dot{F}^{k\ell}g_{k\ell}\right)h_{ij} - 2F\left((h^2)_{ij}+ g_{ij}\right),  \label{s2.evlh2}
		\end{align}

\begin{corollary}\label{lem.cond-ii}
Let $F:\Gamma_+\to (0,\infty)$ be a smooth symmetric function. If $F$ satisfies the condition \eqref{equ-property3}, then 
\begin{equation}\label{eq:old-property2}
\big(\dot F^i\kappa_i-\dot F^j\kappa_j\big)(\kappa_i-\kappa_j)\ge 0
\qquad \text{for all } i\neq j.
\end{equation}
\end{corollary}

\begin{equation}
\big(\dot F^i\kappa_i-\dot F^j\kappa_j\big)(\kappa_i-\kappa_j)\ge 0
\qquad \text{for all } i\neq j.
\end{equation}

\begin{proposition}\label{lem:gm-lower-bound}
 Let $F:\Gamma_+\to (0,\infty)$ be a smooth symmetric function on $\Gamma_+$ satisfying the Assumption \ref{assumption}. Then
\begin{equation}\label{equ-f-lower}
F(\kappa)\ge \left(\prod_{i=1}^n \kappa_i\right)^{1/n}
\qquad \text{for all }\kappa\in \Gamma_+.
\end{equation}
\end{proposition}

\begin{equation}
F(\kappa)\ge \left(\prod_{i=1}^n \kappa_i\right)^{1/n}
\qquad \text{for all }\kappa\in \Gamma_+.
\end{equation}

\begin{proposition}\label{s4.lem-pinc}
	Let $F$ satisfy Assumption \ref{assumption}. If the initial hypersurface $M_0$ has positive sectional curvature, then there exists a constant $C$ depending on $n$ and $M_0$ such that along the flow \eqref{equ-f} in $\mathbb{H}^{n+1}\ (n\geq 3)$ the evolving hypersurface $M_t$ satisfies
	\begin{equation}
		\label{equ-pinching}
		\kappa_n \leq C \kappa_1
	\end{equation}
    for all $t>0$.
\end{proposition}

\begin{proof}
	The sectional curvature defines a smooth function on the Grassmannian bundle of two-dimensional subspaces of the tangent bundle $TM$. For convenience we lift it to a function on the orthonormal frame bundle $O(M)$ over $M$: Given a point $x\in M$, a time $t\geq 0$, and a frame $\mathbb{O}=\{e_1,e_2,\cdots,e_n\}$ for $T_x M$ which is orthonormal with respect to $g(x,t)$, we define
	\begin{equation*}
		G(x,t,\mathbb{O}) = h_{(x,t)}(e_1,e_1)h_{(x,t)}(e_2,e_2) - h_{(x,t)}(e_1,e_2)^2 - 1- \varepsilon F(x,t)^2,
	\end{equation*}
    where $\varepsilon\in(0,\varepsilon_0)$ is small enough such that $G>0$ at $t=0$. Here
    \begin{equation*}
    	\varepsilon_0 = \min\limits_{M_0}\dfrac{\kappa_1\kappa_2-1}{F^2}.
    \end{equation*}
    Consider a point $(x_0,t_0)$ and a frame $\mathbb{O}_0=\{\overline{e}_1,\overline{e}_2,\cdots,\overline{e}_n\}$ at which a new minimum of the function $G$ is attained, so that $$G(x,t,\mathbb{O})\geq G(x_0,t_0,\mathbb{O}_0)$$
    for all $x\in M$, all $t\in [0,t_0]$, and all $\mathbb{O}\in F(M)_{(x,t)}$. Since $\mathbb{O}_0$ achieves the minimum of $G$ over the fiber $F(M)_{(x_0,t_0)}$, the vectors $\overline{e}_1$ and $\overline{e}_2$ can be rotated to become eigenvectors of $h_{(x_0,t_0)}$ corresponding to $\kappa_1$ and $\kappa_2$, where $\kappa_1\leq\cdots\leq \kappa_n$ are the principal curvatures at $(x_0,t_0)$. Moreover, $G$ is invariant under rotation in the subspace orthogonal to $\overline{e}_1$ and $\overline{e}_2$. Consequently, we may assume without loss of generality that $h(\overline{e}_i,\overline{e}_i)=\kappa_i$ and $h(\overline{e}_i,\overline{e}_j)=0$ for $i\neq j$.
    
    We derive the differential inequality satisfied by $G$ at the minimal point at $(x_0,t_0,\mathbb{O}_0)$. Note that by the evolution equation \eqref{s2.evl-g} of the metric, the frame $\mathbb{O}(t)$ for $T_x M$ defined by
    \begin{equation*}
    	\dfrac{\mathrm{d}}{\mathrm{d}t}e_i(t) = F \mathcal{W}(e_i)
    \end{equation*}
    remains orthonormal with respect to $g(x,t)$ if $e_i(t_0)=\overline{e}_i$ for each $i$.  We first compute the 
    time derivative of $G$ at $(x_0,t_0,\mathbb{O}_0)$. Combining \eqref{s2.evlh2} and \eqref{equ-psi}, we have 
    \begin{align}
    	\left.\dfrac{\partial}{\partial t} G \right|_{(x_0,t_0,\mathbb{O}_0)} = & \kappa_1 \left(\frac{\partial}{\partial t}h_{22}+2h\left(\frac{\partial}{\partial t}e_2,e_2\right)\right) \nonumber\\
        &+ \kappa_2 \left(\frac{\partial}{\partial t}h_{11}+2h\left(\frac{\partial}{\partial t}e_1,e_1\right)\right) - 2 \varepsilon F \dfrac{\partial}{\partial t}F \notag\\
    	= & \kappa_1 \left(\frac{\partial}{\partial t}h_{22}+2F\kappa_2^2\right) \nonumber\\
        &+ \kappa_2 \left(\frac{\partial}{\partial t}h_{11}+2F\kappa_1^2\right) - 2 \varepsilon F \dfrac{\partial}{\partial t}F \notag\\
        = &  \kappa_1 \dot{F}^{k\ell} \nabla_k\nabla_{\ell} h_{22} + \kappa_2 \dot{F}^{k\ell}\nabla_k\nabla_{\ell} h_{11} \nonumber\\
        &+ \kappa_1\ddot{F}^{k\ell,pq}\nabla_2 h_{k\ell}\nabla_2 h_{pq} + \kappa_2 \ddot{F}^{k\ell,pq}\nabla_1 h_{k\ell}\nabla_1 h_{pq} \nonumber\\
    	& + 2 \left( \dot{F}^{k\ell}(h^2)_{k\ell} + \dot{F}^{k\ell}g_{k\ell} \right)\kappa_1\kappa_2 - 2 F (\kappa_1+\kappa_2) \nonumber\\
    	& - 2\varepsilon F \left( \dot{F}^{k\ell} \nabla_k \nabla_{\ell} F + \left( \dot{F}^{k\ell}(h^2)_{k\ell} - \dot{F}^{k\ell}g_{k\ell} \right)F  \right).
    	\label{equ-Gt-2}
    \end{align}
    
    The zero-order terms in \eqref{equ-Gt-2} satisfy
    \begin{align}
    	& 2 \left( \dot{F}^{k\ell}(h^2)_{k\ell} + \dot{F}^{k\ell}g_{k\ell} \right)\kappa_1\kappa_2 - 2 F (\kappa_1+\kappa_2) - 2 \varepsilon F^2 \left( \dot{F}^{k\ell}(h^2)_{k\ell} - \dot{F}^{k\ell}g_{k\ell} \right) \nonumber\\
    	=& 2 \sum\limits_{k=1}^n \dot{F}^k \kappa_k^2 G + 2 \sum\limits_{k=1}^n \dot{F}^k (\kappa_k-\kappa_1)(\kappa_k-\kappa_2) + 2 \varepsilon F^2 \dot{F}^{k\ell}g_{k\ell} ~
    	\geq ~ 0.   \label{equ-zero-2}
    \end{align}
    To estimate the Hessian terms in \eqref{equ-Gt-2}, we consider the second derivatives of $G$ along a curve on $O(M)$ defined as follows: We let $\gamma$ be any geodesic of $g(t_0)$ in $M$ with $\gamma(0)=x_0$, and define a frame $\mathbb{O}(s)=(e_1(s),e_2(s),\cdots,e_n(s))$ at $\gamma(s)$ by taking $e_i(0)=\overline{e}_i$ for each $i$, and
    \begin{equation*}
    	\nabla_s e_i(s) = \Gamma_{ij}e_j(s)
    \end{equation*}
    for some constant antisymmetric matrix $\Gamma$. Then we compute
    \begin{align}
    	\left.\dfrac{\mathrm{d}^2}{\mathrm{d}s^2}G(x(s),t_0,\mathbb{O}(s))\right|_{s=0} = & \kappa_2 \nabla_s^2 h_{11} + \kappa_1 \nabla_s^2 h_{22} + 2 \left( \nabla_s h_{22}\nabla_s h_{11} - (\nabla_s h_{12})^2 \right) \nonumber\\
    	& + 4 \sum\limits_{p>2} \Gamma_{1p}\kappa_2 \nabla_s h_{1p} + 4 \sum\limits_{p>2} \Gamma_{2p}\kappa_1 \nabla_s h_{2p} \nonumber\\
    	&+ 2 \sum\limits_{p>2} \Gamma_{1p}^2 \kappa_2(\kappa_p-\kappa_1) + 2 \sum\limits_{p>2} \Gamma_{2p}^2 \kappa_1 (\kappa_p-\kappa_2) \nonumber\\
    	& - 2 \varepsilon F \nabla_s^2 F - 2 \varepsilon (\nabla_s F)^2.  
    	\label{equ-d2G-2}
    \end{align}
    Since $G$ has a minimum at $(x_0,t_0,\mathbb{O}_0)$, the right hand side of \eqref{equ-d2G-2} is nonnegative for any choice of $\Gamma$. The optimal choice of $\Gamma$, which minimizes the expression \eqref{equ-d2G-2}, is given by
    \begin{equation*}
    	\Gamma_{kp}=\begin{cases}
    		-\dfrac{\nabla_s h_{kp}}{\kappa_p-\kappa_k},&k\in\{1,2\},p>2,\\
    		\dfrac{\nabla_s h_{kp}}{\kappa_k-\kappa_p},&p\in\{1,2\},k>2,\\
    		0,& \mbox{otherwise}.
    	\end{cases}
    \end{equation*}
    Thus, minimizing over $\Gamma$ in \eqref{equ-d2G-2} gives
    \begin{align}
    	0 \leq & \kappa_2 \nabla_s^2 h_{11} + \kappa_1 \nabla_s^2 h_{22} + 2 \left( \nabla_s h_{22}\nabla_s h_{11} - (\nabla_s h_{12})^2  \right) \nonumber\\
    	& - 2 \sum\limits_{p>2} \dfrac{\kappa_2}{\kappa_p-\kappa_1}(\nabla_s h_{1p})^2 - 2 \sum\limits_{p>2} \dfrac{\kappa_1}{\kappa_p - \kappa_2} (\nabla_s h_{2p})^2 \nonumber \\
    	& - 2 \varepsilon F \nabla_s^2 F - 2 \varepsilon (\nabla_s F)^2,
    	\label{equ-2nd-2}
    \end{align}
    with the terms on the second line regarded as vanishing if the denominators vanish (since the corresponding component of $\nabla h$ vanishes in that case). Using the formula \eqref{equ-secf}, we rewrite the gradient terms in \eqref{equ-Gt-2} as follows:
    \begin{align}
       \ddot{F}^{k\ell,pq}\nabla_2 h_{k\ell}\nabla_2 h_{pq}= &\sum\limits_{i,j} \ddot{F}^{ij}\nabla_2 h_{ii}\nabla_2 h_{jj} + 2\sum\limits_{i< j} \dfrac{\dot{F}^i -\dot{F}^{j}}{\kappa_i-\kappa_{j}}(\nabla_2 h_{ij})^2, \\
       \ddot{F}^{k\ell,pq}\nabla_1 h_{k\ell}\nabla_1 h_{pq} =&\sum\limits_{i,j} \ddot{F}^{ij}\nabla_1 h_{ii}\nabla_1 h_{jj} + 2\sum\limits_{i< j} \dfrac{\dot{F}^i -\dot{F}^{j}}{\kappa_i-\kappa_{j}} (\nabla_1 h_{ij})^2. \label{equ-HessF}
    \end{align}
    
    Substituting \eqref{equ-zero-2} and \eqref{equ-2nd-2} -- \eqref{equ-HessF} into \eqref{equ-Gt-2}, we obtain
    \begin{align}
    	\dfrac{\partial}{\partial t}G \bigg|_{(x_0,t_0,\mathbb{O}_0)} \geq &  \kappa_2 \left( \sum\limits_{i,j} \ddot{F}^{ij}\nabla_1 h_{ii}\nabla_1 h_{jj} + 2\sum\limits_{i< j} \dfrac{\dot{F}^i -\dot{F}^{j}}{\kappa_i-\kappa_{j}} (\nabla_1 h_{ij})^2 \right) \notag\\
    	& + \kappa_1 \left( \sum\limits_{i,j} \ddot{F}^{ij}\nabla_2 h_{ii}\nabla_2 h_{jj} + 2\sum\limits_{i< j} \dfrac{\dot{F}^i -\dot{F}^{j}}{\kappa_i-\kappa_{j}}(\nabla_2 h_{ij})^2 \right) \notag\\
    	& - 2 \sum\limits_{i=1}^n \dot{F}^i \left( \nabla_i h_{22}\nabla_i h_{11} - (\nabla_i h_{12})^2  \right) \nonumber\\
    	& + 2 \sum\limits_{i=1}^n \dot{F}^i \left( \sum\limits_{p>2} \dfrac{\kappa_2}{\kappa_p-\kappa_1}(\nabla_i h_{1p})^2 + \sum\limits_{p>2} \dfrac{\kappa_1}{\kappa_p-\kappa_2} (\nabla_i h_{2p})^2  \right) \nonumber \\
    	& + 2 \varepsilon \sum\limits_{i=1}^n \dot{F}^i (\nabla_i F)^2. \label{equ-ptg22}
    \end{align}

  Using \eqref{equ-property3} in Assumption \ref{assumption}, we estimate the terms involving second order derivatives of $F$ as follows 
    \begin{align}
        \sum\limits_{i,j} \ddot{F}^{ij}\nabla_1 h_{ii}\nabla_1 h_{jj}\geq &\dfrac{(\nabla_1 F)^2}{F} - \sum\limits_{i=1}^n \dfrac{\dot{F}^i}{\kappa_i}(\nabla_1 h_{ii})^2,\label{s3.d2F}\\
          \sum\limits_{i,j} \ddot{F}^{ij}\nabla_2 h_{ii}\nabla_2 h_{jj}\geq &\dfrac{(\nabla_2 F)^2}{F} - \sum\limits_{i=1}^n \dfrac{\dot{F}^i}{\kappa_i}(\nabla_2 h_{ii})^2. \label{s3.d2Fb}
    \end{align}
Under the Assumption \ref{assumption}, by Corollary \ref{lem.cond-ii} we have
\[
 \big(\dot F^i\kappa_i-\dot F^j\kappa_j\big)(\kappa_i-\kappa_j)\ge 0
\qquad \text{for all } i\neq j.   
\]
This implies that for any $i\neq j$,
    \begin{equation*}
    	\dfrac{\dot{F}^i-\dot{F}^j}{\kappa_i - \kappa_j} + \dfrac{\dot{F}^i}{\kappa_j}=\dfrac{\dot{F}^i\kappa_i-\dot{F}^j\kappa_j}{(\kappa_i-\kappa_j)\kappa_j} \geq 0.%\ \ \mbox{and}\ \ \dfrac{\dot{f}^i-\dot{f}^j}{\kappa_i - \kappa_j} + \dfrac{\dot{f}^j}{\kappa_i} =\dfrac{\dot{f}^i\kappa_i-\dot{f}^j\kappa_j}{(\kappa_i-\kappa_j)\kappa_i}\geq 0.
    \end{equation*}
    Then we estimate
    \begin{align}
    	2 \sum\limits_{i<j} \dfrac{\dot{F}^i - \dot{F}^j}{\kappa_i-\kappa_j}(\nabla_1 h_{ij})^2
    	=& 2\dfrac{\dot{F}^1-\dot{F}^2}{\kappa_1-\kappa_2}(\nabla_2 h_{11})^2 + 2 \sum\limits_{\substack{1\leq i<j\leq n\\j>2}}\dfrac{\dot{F}^i -\dot{F}^j}{\kappa_i -\kappa_j}(\nabla_1 h_{ij})^2 \notag\\
    	\geq& -\left(\dfrac{\dot{F}^1}{\kappa_2}+\dfrac{\dot{F}^2}{\kappa_1}\right)(\nabla_2 h_{11})^2 - 2 \sum\limits_{i=1}^n\sum\limits_{\substack{j>i\\j>2}}\dfrac{\dot{F}^i}{\kappa_j}(\nabla_1 h_{ij})^2.
    	\label{equ-estd1}
    \end{align}
    Similarly,
    \begin{align}
    	2 \sum\limits_{i<j} \dfrac{\dot{F}^i - \dot{F}^j}{\kappa_i-\kappa_j}(\nabla_2 h_{ij})^2 	=& 2\dfrac{\dot{F}^1-\dot{F}^2}{\kappa_1-\kappa_2}(\nabla_1 h_{22})^2 + 2 \sum\limits_{\substack{1\leq i<j\leq n\\j>2}}\dfrac{\dot{F}^i -\dot{F}^j}{\kappa_i -\kappa_j}(\nabla_2 h_{ij})^2 \notag\\
    	\geq& -\left(\dfrac{\dot{F}^1}{\kappa_2}+\dfrac{\dot{F}^2}{\kappa_1}\right)(\nabla_1 h_{22})^2 - 2 \sum\limits_{i=1}^n\sum\limits_{\substack{j>i\\j>2}}\dfrac{\dot{F}^i}{\kappa_j}(\nabla_2 h_{ij})^2.
    	\label{equ-estd2}
    \end{align}
Substituting \eqref{s3.d2F} -- \eqref{equ-estd2} into \eqref{equ-ptg22}, we obtain
    \begin{align}
    	\left.\dfrac{\partial}{\partial t}G\right|_{(x_0,t_0,\mathbb{O}_0)} \geq & \kappa_2 \left( \dfrac{(\nabla_1 F)^2}{F} - \sum\limits_{i=1}^n \dfrac{\dot{F}^i}{\kappa_i}(\nabla_1 h_{ii})^2 \right)\nonumber\\
        &-\kappa_2\left(\left(\dfrac{\dot{F}^1}{\kappa_2}+\dfrac{\dot{F}^2}{\kappa_1}\right)(\nabla_2 h_{11})^2 + 2 \sum\limits_{i=1}^n\sum\limits_{\substack{j>i\\j>2}}\dfrac{\dot{F}^i}{\kappa_j}(\nabla_1 h_{ij})^2  \right) \nonumber\\
    	& +  \kappa_1 \left( \dfrac{(\nabla_2 F)^2}{F} - \sum\limits_{i=1}^n \dfrac{\dot{F}^i}{\kappa_{i}}(\nabla_2 h_{ii})^2\right)\nonumber\\
        &-\kappa_1\left( \left(\dfrac{\dot{F}^1}{\kappa_2}+\dfrac{\dot{F}^2}{\kappa_1}\right)(\nabla_1 h_{22})^2 + 2 \sum\limits_{i=1}^n\sum\limits_{\substack{j>i\\j>2}}\dfrac{\dot{F}^i}{\kappa_j}(\nabla_2 h_{ij})^2 \right) \notag\\
    	& + 2\sum\limits_{i=1}^n \dot{F}^i \left( -\nabla_i h_{22}\nabla_i h_{11} + (\nabla_i h_{12})^2 \right) \notag\\
    	& + 2 \sum\limits_{i=1}^n \dot{F}^{i}\left( \sum\limits_{j>2}\dfrac{\kappa_2}{\kappa_j-\kappa_1}(\nabla_i h_{1j})^2 +\sum\limits_{j>2}\dfrac{\kappa_1}{\kappa_j-\kappa_2}(\nabla_i h_{2j})^2 \right) \notag\\
    	& + 2\varepsilon \sum\limits_{i=1}^n \dot{F}^i (\nabla_i F)^2. \notag\\
    	\geq & \kappa_2\dfrac{(\nabla_1 F)^2}{F} + \kappa_1 \dfrac{(\nabla_2 F)^2}{F} \nonumber\\
        &- \sum\limits_{i=1}^2 \dfrac{\kappa_2}{\kappa_1}\dot{F}^i (\nabla_i h_{11})^2 - \sum\limits_{i=1}^2 \dfrac{\kappa_1}{\kappa_2}\dot{F}^i (\nabla_i h_{22})^2 \nonumber\\
    	& - 2 \sum\limits_{i=1}^n \dot{F}^i \nabla_i h_{22}\nabla_i h_{11} + 2\varepsilon \sum\limits_{i=1}^n \dot{F}^i (\nabla_i F)^2 \nonumber\\
    	= & \kappa_2 \dfrac{(\nabla_1 F)^2}{F} + \kappa_1 \dfrac{(\nabla_2 F)^2}{F} - \sum\limits_{i=1}^2 \dfrac{\dot{F}^i}{\kappa_1 \kappa_2}(\kappa_2 \nabla_i h_{11}+ \kappa_1 \nabla_i h_{22})^2 \nonumber\\
    	& - 2 \sum\limits_{i=3}^n \dot{F}^i \nabla_i h_{22}\nabla_i h_{11} + 2\varepsilon \sum\limits_{i=1}^n \dot{F}^i (\nabla_i F)^2.  \label{equ-2ndgauss2}
    \end{align}
    
    Since $(x_0,\mathbb{O}_0)$ is a minimum point of $G$ at time $t_0$, we have $\nabla_i G=0$ for $i=1,2,\cdots,n$, which yields the critical condition
    \begin{equation}
    	\kappa_2 \nabla_i h_{11} + \kappa_1 \nabla_i h_{22} =2\varepsilon F\nabla_i F,\ \ i=1,\cdots,n.
    	\label{equ-critical-2}
    \end{equation} 
	The fact that $G\geq 0$ at $(x_0,\mathbb{O}_0)$ implies 
    \begin{equation}\label{s4.Ggeq0}
        \kappa_1\kappa_2 \geq 1+\varepsilon F^2\geq \varepsilon F^2
    \end{equation}
holds at the minimum point $x_0$.    Using \eqref{equ-critical-2} and \eqref{s4.Ggeq0},  we estimate the terms in \eqref{equ-2ndgauss2} involving $i=1,2$ as follows:
    \begin{align}
    	&\kappa_2 \dfrac{(\nabla_1 F)^2}{F} + \kappa_1 \dfrac{(\nabla_2 F)^2}{F} - \sum\limits_{i=1}^2 \dfrac{\dot{F}^i}{\kappa_1 \kappa_2}(\kappa_2 \nabla_i h_{11}+ \kappa_1 \nabla_i h_{22})^2 + 2\varepsilon \sum\limits_{i=1}^2 \dot{F}^i (\nabla_i F)^2 \nonumber\\
    	\geq & \sum\limits_{i=1}^2 \dfrac{\varepsilon F^2}{\kappa_i}\dfrac{(\nabla_i F)^2}{F} - \dfrac{4\varepsilon^2 F^2}{\kappa_1 \kappa_2} \sum\limits_{i=1}^2 \dot{F}^i (\nabla_i F)^2 +  2\varepsilon \sum\limits_{i=1}^2 \dot{F}^i (\nabla_i F)^2 \nonumber\\
    	\geq & \sum\limits_{i=1}^2 \dfrac{\varepsilon F^2}{\kappa_i}\dfrac{(\nabla_i F)^2}{F} - 2\varepsilon \sum\limits_{i=1}^2 \dot{F}^i (\nabla_i F)^2 \notag\\
    	= & \sum\limits_{i=1}^2 \frac{\varepsilon}{\kappa_i}\left(F-2\dot{F}^{i}\kappa_i\right)(\nabla_i F)^2 \geq 0,
    	\label{equ-est1}
    \end{align}
    where in the last line we used the inequality
    \begin{equation*}
    	F = \sum\limits_{i=1}^n \dot{F}^i \kappa_i \geq 2 \dot{F}^j\kappa_j,\ \ j=1,2,
    \end{equation*}
	which follows from \eqref{eq:old-property2} in Corollary \ref{lem.cond-ii} and $n\geq 3$. For the remaining terms in \eqref{equ-2ndgauss2} with $i\geq 3$, we again employ the critical condition \eqref{equ-critical-2} and \eqref{s4.Ggeq0}.  For each such $i$ we have
    \begin{align*}
    	& -\nabla_i h_{22}\nabla_i h_{11} + \varepsilon (\nabla_i F)^2 \nonumber\\
    	=& \dfrac{\kappa_2}{\kappa_1}(\nabla_i h_{11})^2 - 2\varepsilon\dfrac{F}{\kappa_1}\nabla_i F \nabla_i h_{11} + \varepsilon(\nabla_i F)^2 \notag\\
    	\geq & \dfrac{\varepsilon F^2}{\kappa_1^2}(\nabla_i h_{11})^2 - 2\varepsilon\dfrac{F}{\kappa_1}\nabla_i F \nabla_i h_{11} + \varepsilon(\nabla_i F)^2\notag\\
    	=& \varepsilon \left( \dfrac{F}{\kappa_1}\nabla_i h_{11} - \nabla_i F \right)^2\notag~
    	\geq  0.%\label{equ-criuse}
    \end{align*}
   Consequently,
    \begin{equation}
    	-2\sum\limits_{i=3}^n \dot{F}^i \nabla_i h_{22}\nabla_i h_{11} + 2\varepsilon \sum\limits_{i=3}^n \dot{F}^i (\nabla_i F)^2 \geq 0.
    	\label{equ-est2}
    \end{equation}

    Substituting \eqref{equ-est1} and \eqref{equ-est2} into \eqref{equ-2ndgauss2}, we conclude
    \begin{equation*}
    	\left.\frac{\partial}{\partial t}G\right|_{(x_0,t_0,\mathbb{O}_0)}\geq 0.
    \end{equation*}
    By the maximum principle, the minimum of $G$ cannot decrease. Since initially $G>0$, it remains positive for all $t>0$. Then using the property \eqref{equ-f-lower} in Proposition \ref{lem:gm-lower-bound} for the function $F$,  $M_t$ satisfies
    \begin{equation*}
    	\kappa_1 \kappa_2 >1+ \varepsilon F^2 \geq \varepsilon \left(\prod_{i=1}^n \kappa_i\right)^{\frac{2}{n}} \geq \varepsilon \kappa_1^{\frac{n-2}{n}} \kappa_2 \kappa_n^{\frac{2}{n}}.
    \end{equation*}
  From this, we obtain the desired pinching estimate 
    \[\kappa_n\leq C\kappa_1\] 
    with $C=\varepsilon^{-n/2}$, which depends only on $n$ and $M_0$.
\end{proof}

\end{document}
